% ============ 第1章 §1.3 算子的范数 (原书 pp.9-14) ============
\section{算子的范数}\label{sec:1.3}
\index{J@基本向量}

\index{S@算子范数}






作为线性代数的重要要素, 可以举出线性空间与线性映射. 作为线性空间, 这里考虑了有限集合 $\{1, 2, \ldots, N\}$ 上的 $\ell^{p}$ 空间, $1 \leq p \leq \infty$. 由 $\ell^{p}$-范数, 向量的大小被确定了. 与之对应, 对未必相等的自然数 $M$, $N$ 与满足 $1 \leq p, q \leq \infty$ 的未必相等的实数 $p$, $q$, 想确定从有限集合 $\{1, 2, \ldots, M\}$ 上的 $\ell^{p}$ 空间到有限集合 $\{1, 2, \ldots, N\}$ 上的 $\ell^{q}$ 空间的线性映射 $T$ 的大小. 考虑到 $T$ 的线性, 具体按如下方式确定.
\[
\| T \|_{\ell^{p} \to \ell^{q}} = \sup_{a \in \ell^{p} \setminus \{0\}} \frac{\| T a \|_{\ell^{q}}}{\| a \|_{\ell^{p}}} = \sup_{a \in \ell^{p}, \| a \|_{\ell^{p}} = 1} \| T a \|_{\ell^{q}}.
\]
将 $\|T\|_{\ell^{p} \to \ell^{q}}$ 称为 $T$ 的从 $(\R^{N}, \ell^{p})$ 到 $(\R^{M}, \ell^{q})$ 的算子范数, 但除去根据上下文显然的要素, 例如也称为从 $\ell^{p}$ 到 $\ell^{q}$ 的算子范数.

利用三角函数, 可以计算配备 $\ell^{2}$ 范数的坐标平面 $\R^{2}$ 之间的线性算子的算子范数.

\begin{example}\label{exa:1.3}
设 $\R^{2}$ 配备范数 $\|(x, y)\| = \sqrt{x^{2} + y^{2}}$. 线性映射 $f: \R^{2} \to \R^{2}$ 由 $f(x, y) = (6x - 2y, 8x - 4y)$ 给出. 来计算 $f$ 的算子范数. 该值作为 $\sqrt{x^{2} + y^{2}} = 1$ 时 $\sqrt{(6x-2y)^{2} + (8x-4y)^{2}}$ 的上限给出. 由于 $(6x-2y)^{2} + (8x-4y)^{2} = 100x^{2} - 88xy + 20y^{2}$, 将约束条件 $x^{2} + y^{2} = 1$ 表示为 $x = \cos\theta$, $y = \sin\theta$ 时, 该式的值可用适当的 $\alpha$ 合成为
\begin{align*}
100\cos^{2}\theta - 88\sin\theta\cos\theta + 20\sin^{2}\theta & = 60 + 40\cos 2\theta - 44\sin 2\theta \\
& = 60 + 4\sqrt{221}\sin(2\theta + \alpha).
\end{align*}
因此, $f$ 的算子范数为 $2\sqrt{15 + \sqrt{221}} = \sqrt{26} + \sqrt{34}$.
\end{example}

来确认: 即使是相同的线性算子, 若配备的范数不同, 算子范数也不同.

\begin{example}\label{exa:1.4}
线性映射 $f: \R^{2} \to \R^{2}$ 由 $f(x, y) = (4x - 2y, x - 4y)$ 给出.

(1) 在条件 $x^{2} + y^{2} = 1$ 下求 $(4x-2y)^{2} + (x-4y)^{2}$ 的最大值, 为 $\frac{37 + 3\sqrt{65}}{2}$, 所以若定义域与值域都配备 $\ell^{2}$-范数, 则 $f$ 的算子范数为 $\sqrt{\frac{37 + 3\sqrt{65}}{2}}$.

(2) 在条件 $|x| + |y| = 1$ 下求 $(4x-2y)^{2} + (x-4y)^{2}$ 的最大值, 为 $20$, 所以若定义域配备 $\ell^{1}$-范数, 值域配备 $\ell^{2}$-范数, 则 $f$ 的算子范数为 $2\sqrt{5}$.

(3) 在条件 $x^{2} + y^{2} = 1$ 下求 $|4x - 2y| + |x - 4y|$ 的最大值. 令 $k = |4x - 2y| + |x - 4y|$. 设 $k > 0$. 此时, $xy$ 平面上的方程 $k = |4x - 2y| + |x - 4y|$ 是以 $\pm\left(\frac{2k}{7}, \frac{k}{14}\right)$, $\pm\left(\frac{k}{7}, \frac{2k}{7}\right)$ 为顶点的平行四边形. 连接 $\left(\frac{2k}{7}, \frac{k}{14}\right)$ 与 $\left(\frac{k}{7}, \frac{2k}{7}\right)$ 的直线 $k = 3x + 2y$ 到原点的距离为 $\frac{k}{\sqrt{13}}$, 连接 $\left(\frac{2k}{7}, \frac{k}{14}\right)$ 与 $-\left(\frac{k}{7}, \frac{2k}{7}\right)$ 的直线 $k = 5x - 6y$ 到原点的距离为 $\frac{k}{\sqrt{61}}$. 而且, 给出后者的点位于 $\left(\frac{2k}{7}, \frac{k}{14}\right)$ 与 $-\left(\frac{k}{7}, \frac{2k}{7}\right)$ 这两点之间. 因此, 使 $k = |4x - 2y| + |x - 4y|$ 与 $x^{2} + y^{2} = 1$ 相交的 $k$ 的最大值为 $\sqrt{61}$. 所以, 若定义域配备 $\ell^{2}$-范数, 值域配备 $\ell^{1}$-范数, 则 $f$ 的算子范数为 $\sqrt{61}$.

(4) 在条件 $x^{2} + y^{2} = 1$ 下求 $\max(|4x-2y|, |x-4y|)$ 的最大值, 为 $2\sqrt{5}$, 所以若定义域配备 $\ell^{2}$-范数, 值域配备 $\ell^{\infty}$-范数, 则 $f$ 的算子范数为 $2\sqrt{5}$.
\end{example}

\begin{example}\label{exa:1.5}
定义域与值域都配备 $\ell^{2}$-范数. 求由
\[
(a_{1}, a_{2}, a_{3}, a_{4}) \mapsto \left(\frac{a_{1} + a_{2}}{2}, \frac{a_{1} + a_{2} + a_{3} + a_{4}}{4}, 0, 0\right)
\]
给出的从 $\R^{4}$ 到 $\R^{4}$ 的线性映射的范数.

在条件
\[
a_{1}^{2} + a_{2}^{2} + a_{3}^{2} + a_{4}^{2} = 1
\]
下, 所求范数是
\[
\left(\frac{a_{1} + a_{2}}{2}\right)^{2} + \left(\frac{a_{1} + a_{2} + a_{3} + a_{4}}{4}\right)^{2}
\]
的最大值的平方根. 令 $a_{1} = r\cos\theta$, $a_{2} = r\sin\theta$, $a_{3} = R\cos\phi$, $a_{4} = R\sin\phi$, 则条件可改写为 $r^{2} + R^{2} = 1$, 所求范数的平方是
\[
\left(\frac{r\cos\theta + r\sin\theta}{2}\right)^{2} + \left(\frac{r\cos\theta + r\sin\theta + R\cos\phi + R\sin\phi}{4}\right)^{2}
\]
关于 $R$, $r$, $\theta$, $\phi$ 的最大值. 关于这个量, 由三角函数的合成可知,
\[
\frac{r^{2}}{2} + \frac{(r + R)^{2}}{8}
\]
是上界. 于是, 令 $r = \cos\tau$, $R = \sin\tau$, 则
\[
\frac{r^{2}}{2} + \frac{(r + R)^{2}}{8} = \frac{5}{8}\cos^{2}\tau + \frac{1}{4}\sin\tau\cos\tau + \frac{1}{8}\sin^{2}\tau.
\]
由半角及倍角公式,
\[
\frac{r^{2}}{2} + \frac{(r + R)^{2}}{8} = \frac{3}{8} + \frac{1}{8}\sin 2\tau + \frac{1}{4}\cos 2\tau,
\]
所以所求最大值为
\[
\sqrt{\frac{3 + \sqrt{5}}{8}} = \frac{1 + \sqrt{5}}{4}.
\]
\end{example}

用 $\mathrm{diag}(\alpha_{1}, \alpha_{2}, \ldots, \alpha_{N})$ 表示对角成分为 $\alpha_{1}, \alpha_{2}, \ldots, \alpha_{N}$ 的对角矩阵. 来计算对角矩阵生成的线性算子的算子范数.

\begin{example}\label{exa:1.6}
设 $1 \leq p \leq \infty$, $b = (b_{1}, b_{2}, \ldots, b_{N}) \in \R^{N}$ 为常数. 考虑把 $a = (a_{1}, a_{2}, \ldots, a_{N}) \in \R^{N}$ 移到 $m_{b}(a) = (b_{1}a_{1}, b_{2}a_{2}, \ldots, b_{N}a_{N}) \in \R^{N}$ 的映射 $m_{b}$. 用适当的 $1 \leq K \leq N$, 记 $\max_{1 \leq j \leq N} |b_{j}| = |b_{K}|$. 只有第 $K$ 个成分为 $1$ 而其他成分为 $0$ 的数列 (第 $K$ 基本向量) $\mathbf{e}_{K}$ 被映射到 $b_{K}\mathbf{e}_{K}$, 所以 $\|m_{b}\|_{\ell^{p} \to \ell^{p}} \geq |b_{K}|$. 另一方面, 将 $|b_{1}|, |b_{2}|, \ldots, |b_{N}|$ 替换为 $|b_{K}|$ 来考虑, 得 $\|m_{b}(a)\|_{\ell^{p}} \leq |b_{K}| \|a\|_{\ell^{p}}$. 因此, $\|m_{b}\|_{\ell^{p} \to \ell^{p}} = |b_{K}| = \max_{1 \leq j \leq N} |b_{j}| = \|b\|_{\ell^{\infty}}$.
\end{example}

在考虑下一个例子时, 将 $N$ 阶单位矩阵记为 $E$, 来复习一下矩阵的知识. 正交矩阵是指由 $PP^{t} = P^{t}P = E$ 给出的实矩阵. 对所有数列 (向量) $a$, 正交矩阵 $P$ 满足 $\|Pa\|_{\ell^{2}} = \|a\|_{\ell^{2}}$. 对称矩阵是指满足 $A = A^{t}$ 的实矩阵.

\begin{example}\label{exa:1.7}
设 $A$ 为 $N$ 阶实对称矩阵, 则用正交矩阵 $P$ 可将其对角化为 $P^{-1}AP = \mathrm{diag}(\alpha_{1}, \alpha_{2}, \ldots, \alpha_{N})$. 其中设 $|\alpha_{1}| \geq |\alpha_{2}| \geq \cdots \geq |\alpha_{N}|$. 设 $A$, $P^{-1}$, $P^{-1}AP$ 给出的 $\R^{N}$ 上的线性映射分别为 $f$, $f_{P^{-1}}$, $f_{P^{-1}AP}$, 则
\[
\frac{\| f(a) \|_{\ell^{2}}}{\| a \|_{\ell^{2}}} = \frac{\| f_{P^{-1}AP}(f_{P^{-1}}(a)) \|_{\ell^{2}}}{\| f_{P^{-1}}(a) \|_{\ell^{2}}},
\]
所以由例~\ref{exa:1.6} 求得 $\|f\|_{\ell^{2} \to \ell^{2}} = |\alpha_{1}|$.
\end{example}

\begin{exercise}\label{exer:1.10}
在 $\R^{2}$ 上配备范数 $\|(x, y)\| = \sqrt{x^{2} + y^{2}}$. 计算下列算子的范数.

(1) $(x, y) \in \R^{2} \mapsto (y, x) \in \R^{2}.$

(2) $(x, y) \in \R^{2} \mapsto (2x, 2y) \in \R^{2}.$

(3) $(x, y) \in \R^{2} \mapsto (x + y, \sqrt{2}x) \in \R^{2}.$
\end{exercise}

\begin{exercise}\label{exer:1.11}
设 $\R^{2}$ 为配备 $\|(x, y)\| = \sqrt{x^{2} + y^{2}}$ 的 Euclid 空间. 线性映射 $f: \R^{2} \to \R^{2}$ 由 $f(x, y) = (4x - 2y, 8x - 4y)$ 给出.

(1) 计算 $f$ 的算子范数 $\|f\|$.

(2) 计算 $f^{2}$ 的算子范数 $\|f^{2}\|$.

(3) 考察 $\|f^{2}\|$ 与 $\|f\|^{2}$ 的大小关系.
\end{exercise}

\begin{exercise}\label{exer:1.12}
线性映射 $f: \R^{2} \to \R^{2}$ 由 $f(x, y) = (4x - 3y, x - 4y)$ 给出. 求下列情形下 $f$ 的算子范数.

(1) 定义域与值域都配备 $\ell^{2}$-范数的情形.

(2) 定义域与值域都配备 $\ell^{1}$-范数的情形.

(3) 定义域配备 $\ell^{2}$-范数, 值域配备 $\ell^{\infty}$-范数的情形.

(4) 定义域配备 $\ell^{\infty}$-范数, 值域配备 $\ell^{2}$-范数的情形.
\end{exercise}

\begin{exercise}\label{exer:1.13}
求配备 $\ell^{2}$-范数的 $\R^{2}$ 之间的映射 $f(x, y) = (2x - 3y, 3x + 4y)$ 的算子范数. 若定义域配备 $\ell^{\infty}$-范数, 值域配备 $\ell^{2}$-范数, 则 $f$ 的算子范数是多少?
\end{exercise}

\begin{exercise}\label{exer:1.14}
求从配备 $\ell^{2}$-范数的 $\R^{2}$ 到配备 $\ell^{2}$-范数的 $\R^{3}$ 的映射 $f(x, y) = (2x + y, x - 2y, x + y)$ 的算子范数.
\end{exercise}

\begin{exercise}\label{exer:1.15}
求从配备 $\ell^{3}$-范数的 $\R^{2}$ 到 $\R$ 的映射 $f(x, y) = 2x + y$ 的算子范数.
\end{exercise}

\begin{exercise}\label{exer:1.16}
求函数 $f(x, y) = \sqrt[3]{|x+1|^{3} + |y-4|^{3}} - \sqrt[3]{|x+2|^{3} + |y-3|^{3}}$, $x, y \in \R$ 的最大值.
\end{exercise}
